\chapter{Deret Fourier}\label{ch:b2:fourier}

Dapatkah setiap sinyal berkala \omterm{def:b2:structures:generated}{dibangun} ulang dari sinus dan kosinus murni?
Jawaban ``ya'' Fourier yang berani menciptakan satu abad analisis. Bab
ini membuktikan dua pilarnya yang terjangkau pada tingkat ini:
\emph{teorema Dirichlet} (yakni rekonstruksi \omterm{def:b2:funcseq:def}{titik demi titik} untuk fungsi
$C^1$ sepotong-sepotong, lewat \omterm{lem:b2:fourier:kernel}{kernel Dirichlet}) dan \emph{kesamaan Parseval}
(yakni energi sebuah sinyal sama dengan jumlah energi harmoniknya), lalu
memanen deret numerik klasiknya --- dengan Basel $\sum 1/n^2 =
\pi^2/6$ yang pertama di antaranya.

Di sepanjang bab ini, fungsinya berkala-$2\pi$, \omterm{def:b2:metric:continuity}{kontinu} sepotong-sepotong,
dan bernilai kompleks; sedangkan $\mathcal{C}$ menyatakan yang \omterm{def:b2:metric:continuity}{kontinu}.

\section{Koefisien Fourier}

\begin{definition}\label{def:b2:fourier:coefficients}
Adapun \emph{koefisien Fourier}\index{koefisien Fourier} bagi $f$
adalah
\[
c_n(f) = \frac{1}{2\pi}\int_{-\pi}^{\pi} f(t)\,\eu^{-\iu n t}\,\dd
t \qquad (n \in \Z),
\]
dan koefisien bentuk realnya $a_n = c_n + c_{-n}$, $b_n = \iu(c_n
- c_{-n})$, sehingga \emph{jumlah parsial Fourier}\index{jumlah parsial
Fourier} berbunyi
\[
S_N(f)(t) = \sum_{n=-N}^{N} c_n(f)\,\eu^{\iu nt}
= \frac{a_0}{2} + \sum_{n=1}^{N}\bigl(a_n\cos nt + b_n \sin
nt\bigr).
\]
Pada $\mathcal{C}$, definisikan \omterm{def:b2:hermitian:def}{hasil kali dalam Hermitian} $\langle f,
g\rangle = \frac{1}{2\pi}\int_{-\pi}^{\pi}\conj f\,g$: maka
eksponensial $e_n(t) = \eu^{\iu nt}$ bersifat \emph{ortonormal}
(sebab $\langle e_m, e_n\rangle = \delta_{mn}$ lewat perhitungan langsung), dan
$c_n(f) = \langle e_n, f\rangle$: jadi analisis Fourier adalah geometri
\omterm{def:b2:hermitian:adjoint}{Hermitian} (\cref{ch:b2:hermitian}) dalam dimensi tak hingga.
\end{definition}

\begin{proposition}[Ketaksamaan Bessel]\label{prop:b2:fourier:bessel}
Jumlah $S_N(f)$ adalah proyeksi ortogonal $f$ ke ruang
$\mathcal{T}_N$ berisi polinomial trigonometri berderajat $\leq N$,
dan\index{ketaksamaan Bessel}
\[
\sum_{n=-N}^{N} \abs{c_n(f)}^2 \leq \norm f_2^2 =
\frac{1}{2\pi}\int_{-\pi}^{\pi} \abs f^2 :
\]
jadi deret $\sum \abs{c_n}^2$ konvergen, dan $c_n(f) \to 0$ ketika
$\abs n \to \infty$ (yakni Riemann--Lebesgue bagi koefisiennya).
\end{proposition}

\begin{proof}
Selisih $f - S_N(f)$ ortogonal terhadap tiap $e_k$ dengan $\abs k \leq N$ (sebab $\langle
e_k, f - S_N f\rangle = c_k - c_k = 0$): jadi $S_Nf$ adalah proyeksi
ortogonal ke $\mathcal{T}_N = \operatorname{Vect}(e_{-N}, \dots,
e_N)$ (lewat teorema proyeksi jilid Tahun ke-1, kata demi kata pada
latar \omterm{def:b2:hermitian:adjoint}{Hermitiannya}). Menurut Pythagoras: $\norm f_2^2 = \norm{S_Nf}_2^2 +
\norm{f - S_Nf}_2^2 \geq \norm{S_Nf}_2^2 = \sum_{\abs n \leq N}
\abs{c_n}^2$; lalu lewatkan $N \to \infty$.
\end{proof}

\begin{example}[Hampiran terbaik, ditakar]\label{ex:b2:fourier:bestapprox}
Seberapa baik polinomial trigonometri berderajat rendah menghampiri
gigi gergaji $f(t) = t$ (pada $\intoo{-\pi}{\pi}$) dalam
rerata kuadratiknya? Menurut \cref{prop:b2:fourier:bessel}, hampiran
terbaik berderajat $N$ \emph{adalah} $S_N(f)$, dengan galat kuadrat
\[
\norm{f - S_Nf}_2^2 = \norm f_2^2 -
\sum_{\abs n\leq N}\abs{c_n}^2 .
\]
Di sini $\norm f_2^2 = \frac{1}{2\pi}\int_{-\pi}^\pi t^2\dd t =
\frac{\pi^2}{3}$, dan dari $b_n = \frac{2(-1)^{n+1}}{n}$
(\cref{ex:b2:fourier:basel}) diperoleh $\abs{c_n}^2 + \abs{c_{-n}}^2 =
\frac{b_n^2}{2} = \frac{2}{n^2}$. Jadi
\[
\norm{f - S_Nf}_2^2
= \frac{\pi^2}{3} - \sum_{n=1}^{N}\frac{2}{n^2}
\qquad\text{: secara numerik } 1.29,\ 0.79,\ 0.57,\ 0.44
\]
untuk $N = 1, 2, 3, 4$ --- yang menurun, tetapi lambat: sebab ekornya
$\sum_{n>N}\frac2{n^2} \sim \frac2N$ dikendalikan oleh peluruhan
$\frac1n$ yang lambat pada koefisiennya, dan itu sendiri tanda
lompatannya (\cref{exo:b2:fourier:6} yang dibaca terbalik). Pelajaran
penutupnya: Parseval mengubah mutu hampiran menjadi ekor sebuah
deret numerik --- lalu meramalkan, sebelum sebuah gambar pun, bahwa
lompatan membuat deret Fourier konvergen dengan enggan.
\end{example}

\begin{method}[Menghitung koefisien Fourier secara efisien]\label{met:b2:fourier:coefficients}
Sebelum mengintegralkan apa pun:
\begin{enumerate}
  \item \emph{Keparitasan:} $f$ yang genap punya $b_n = 0$, dan $f$ yang ganjil punya
        $a_n = 0$, dan integral yang bertahan menyusut menjadi $\frac2\pi
        \int_0^\pi$ --- jadi separuh kerja dengan keandalan dua kali lipat.
  \item \emph{Polinomial trigonometri sudah selesai:}
        linearkan hasil kalinya ($\cos^3$, $\sin^2\cos$, \dots) lalu
        bacalah koefisiennya
        (\cref{exo:b2:fourier:9}); sebab keortonormalannya membuat
        pengintegralan lanjutan menjadi mubazir.
  \item \emph{Eksponensial kompleks untuk eksponensial:} untuk
        faktor $\eu^{at}$ atau osilasi teredam, hitunglah
        $c_n$ secara langsung --- sebab satu integral $\eu^{(a - \iu
        n)t}$ mengalahkan dua kali pengintegralan parsial
        (\cref{exo:b2:fourier:10}).
  \item \emph{Turunkan sebuah uraian yang diketahui:} bila $f'$ punya
        koefisien yang diketahui dan $f$ \omterm{def:b2:metric:continuity}{kontinu}, maka $c_n(f) =
        \frac{c_n(f')}{\iu n}$ (untuk $n \neq 0$) memulihkan semuanya kecuali
        $c_0$, yang tak lain rerataannya --- ini kerap rute tercepat,
        dan sah persis di bawah hipotesis
        \cref{thm:b2:fourier:parseval}~(1).
\end{enumerate}
\end{method}

\section{Teorema Dirichlet}

\begin{lemma}[Kernel Dirichlet]\label{lem:b2:fourier:kernel}
Berlaku $S_N(f)(x) = \frac{1}{2\pi}\int_{-\pi}^{\pi} f(x + u)\,D_N(u)\,\dd
u$, dengan\index{kernel Dirichlet}
\[
D_N(u) = \sum_{n=-N}^{N} \eu^{\iu nu}
= \frac{\sin\bigl((N + \frac12)u\bigr)}{\sin\frac u2}
\quad (u \notin 2\pi\Z),
\qquad
\frac{1}{2\pi}\int_{-\pi}^{\pi} D_N = 1 .
\]
\end{lemma}

\begin{proof}
Sisipkan definisi $c_n$ ke dalam $S_N$ lalu tukarkan jumlah dan integralnya
(yang sah sebab jumlahnya berhingga):
\[
S_N(f)(x)
= \sum_{n=-N}^{N}\Bigl(\frac{1}{2\pi}\int_{-\pi}^{\pi}
f(t)\,\eu^{-\iu nt}\dd t\Bigr)\eu^{\iu nx}
= \frac{1}{2\pi}\int_{-\pi}^{\pi} f(t)
\sum_{n=-N}^{N}\eu^{\iu n(x-t)}\,\dd t ;
\]
sulihkan $u = t - x$ lalu geser kembali ruas pengintegralannya
ke $\intcc{-\pi}{\pi}$ berkat kekalaan-$2\pi$ integrannya;
sedangkan jangkauan indeksnya yang simetris membuat $\sum_n\eu^{-\iu nu} = D_N(u)$.
Bentuk tertutupnya: yakni jumlah geometri berasio $\eu^{\iu u}$,
\[
D_N(u) = \eu^{-\iu Nu}\,\frac{\eu^{\iu(2N+1)u} - 1}{\eu^{\iu u} -
1}
= \frac{\eu^{\iu(N + \frac12)u} - \eu^{-\iu(N+\frac12)u}}
{\eu^{\iu u/2} - \eu^{-\iu u/2}} ,
\]
yang tak lain hasil bagi sinusnya. Rerataannya $1$: sebab hanya $n = 0$
yang menyumbang.
\end{proof}

\begin{theorem}[Lema Riemann--Lebesgue]\label{thm:b2:fourier:riemannlebesgue}
Untuk $g$ yang \omterm{def:b2:metric:continuity}{kontinu} sepotong-sepotong pada sebuah ruas, berlaku
$\int_a^b g(t)\sin(\lambda t + \varphi)\,\dd t \to 0$ ketika $\lambda
\to +\infty$.\index{lema Riemann--Lebesgue}
\end{theorem}

\begin{proof}
Hampiri $g$ secara \omterm{def:b2:funcseq:def}{seragam} dengan fungsi tangga
(sebab \cref{thm:b2:funcseq:weierstrass} tak diperlukan --- hampiran
fungsi tangga yang dasar bagi fungsi \omterm{def:b2:metric:continuity}{kontinu} sepotong-sepotong
sudah cukup) lalu integralkan tiap tangganya secara gamblang: maka setiap
kepingnya menyumbang $O\bigl(\frac1\lambda\bigr)$, dan galat
hampirannya menyumbang $\varepsilon(b - a)$. Hujah ini dikerjakan
selengkapnya sebagai latihan terakhir bab integrasi pada jilid
Tahun ke-1; sedangkan untuk keping $C^1$ kita malah dapat mengintegralkan
secara parsial lalu membatasinya dengan $\frac C\lambda$.
\end{proof}

\begin{example}[Secepat apa koefisiennya mati?]\label{ex:b2:fourier:decaytable}
Riemann--Lebesgue mengatakan koefisiennya menuju $0$; sedangkan
\emph{lajunya} merupakan alat ukur kemulusan. Tiga spesimen dari bab
ini dan latihannya:
\[
\text{gelombang persegi: } b_n = \frac{4}{\pi n}\ (n\ \text{ganjil}),
\qquad
\abs t : a_n = \frac{-4}{\pi n^2}\ (n\ \text{ganjil}),
\qquad
\abs{\sin t} : a_{2k} = \frac{-4}{\pi(4k^2-1)} .
\]
Sebuah lompatan pada $f$ (gelombang persegi, gigi gergaji) meninggalkan
koefisien berorde $\frac1n$: jadi tanpa \omterm{def:b2:funcseq:series}{kekonvergenan normal}, dan dengan
lonjakan Gibbs di lompatannya. Sedangkan kekontinuan dengan sebuah sudut ---
yakni lompatan pada $f'$ belaka --- memperbaiki ordenya menjadi
$\frac{1}{n^2}$: sehingga \omterm{def:b2:funcseq:series}{kekonvergenan normal} dan rekonstruksi \omterm{def:b2:funcseq:def}{seragam}.
Secara umum $k$ turunan membeli $c_n =
O(n^{-k})$ (\cref{exo:b2:fourier:6}), dan sebaliknya \omterm{def:b2:reduction:eigen}{spektrum} yang
meluruh lebih cepat daripada setiap pangkat memaksa $f$ menjadi $C^\infty$
(turunkan suku demi suku, kini secara sah). Pelajaran
penutupnya: keteraturan sinyalnya dan peluruhan \omterm{def:b2:reduction:eigen}{spektrumnya}
adalah informasi yang sama --- jadi seorang insinyur membaca yang satu dari
kemiringan yang lain tanpa pernah menggambar fungsinya.
\end{example}

\begin{theorem}[Dirichlet]\label{thm:b2:fourier:dirichlet}
Misalkan $f$ berkala-$2\pi$ dan $C^1$ sepotong-sepotong. Maka untuk setiap $x$,
\[
S_N(f)(x) \xrightarrow[N \to \infty]{}
\frac{f(x^+) + f(x^-)}{2}
\]
(yakni rerata limit sepihaknya) --- khususnya $S_N(f)(x) \to
f(x)$ di setiap titik kekontinuannya.\index{teorema Dirichlet}
\end{theorem}


\begin{proof}
Menurut lema kernelnya dan rerataan satuannya, setelah integralnya dipecah
menjadi paruh $u > 0$ dan $u < 0$ (yang masing-masing bererata $\frac12$):
\begin{align*}
S_N(f)(x) - \frac{f(x^+) + f(x^-)}{2}
&= \frac{1}{2\pi}\int_{0}^{\pi} \bigl(f(x+u) -
f(x^+)\bigr)D_N(u)\,\dd u\\
&\quad+ \frac{1}{2\pi}\int_{-\pi}^{0}\bigl(f(x+u) -
f(x^-)\bigr)D_N(u)\,\dd u .
\end{align*}
Tanganilah yang pertama (sebab yang kedua simetris). Tulislah
\[
\bigl(f(x + u) - f(x^+)\bigr)\,D_N(u)
= \underbrace{\frac{f(x+u) - f(x^+)}{\sin\frac u2}}_{g(u)}\,
\sin\Bigl(\Bigl(N + \frac12\Bigr)u\Bigr) .
\]
Fungsi $g$ \omterm{def:b2:metric:continuity}{kontinu} sepotong-sepotong pada $\intoc{0}{\pi}$ dan
punya \emph{limit hingga di $0^+$}: sebab dengan menulis
\[
g(u) = \frac{f(x+u) - f(x^+)}{u}\cdot\frac{u}{\sin\frac u2} ,
\]
faktor pertamanya menuju $f'(x^+)$ (berkat keterdiferensialan
sepihaknya, dari $C^1$ sepotong-sepotong) dan yang kedua menuju $2$
(lewat limit baku $\frac{\sin v}{v} \to 1$ di $v = \frac
u2$): jadi $g(0^+) = 2f'(x^+)$ ada. Karena itu $g$ meluas secara
\omterm{def:b2:metric:continuity}{kontinu} sepotong-sepotong ke $\intcc{0}{\pi}$, lalu Riemann--Lebesgue
(\cref{thm:b2:fourier:riemannlebesgue}) mengirim integralnya ke
$0$. Inilah seluruh maksud hipotesisnya: sebab tanpa
turunan sepihaknya, faktor $\frac{1}{\sin(u/2)}$ meledak
di $0$ lebih cepat daripada yang dapat diimbangi Riemann--Lebesgue, dan
\omterm{def:b2:funcseq:def}{kekonvergenan titik demi titik} sungguh dapat gagal untuk $f$ yang sekadar
\omterm{def:b2:metric:continuity}{kontinu} --- yakni celah yang ditutup teorema Fejér (soal akhir pekan)
lewat perata-rataan.
\end{proof}

\begin{example}[Dirichlet di sebuah lompatan]\label{ex:b2:fourier:jumpvalue}
Untuk gigi gergaji $f(t) = t$ pada $\intoo{-\pi}{\pi}$
(\cref{ex:b2:fourier:basel} di bawah), perluasan berkalanya
melompat di $t = \pi$ dari $f(\pi^-) = \pi$ ke $f(\pi^+) = -\pi$.
Dirichlet menjanjikan nilai $\frac{\pi + (-\pi)}{2} = 0$
di sana, dan memang setiap suku $\sum
\frac{2(-1)^{n+1}}{n}\sin nt$ lenyap di $t = \pi$: jadi
deretnya konvergen dengan sopan ke titik tengahnya, sambil mengabaikan kedua
nilai sepihaknya. Sedangkan memindahkan titik penilaiannya ke $t =
\frac\pi2$ (yakni titik kekontinuan) mengubah deret yang sama
menjadi $\frac\pi4$ milik Leibniz. Satu deret, dua kelakuan ---
persis kedua klausa teoremanya.
\end{example}

\begin{theorem}[Kekonvergenan normal untuk $C^1$; Parseval]\label{thm:b2:fourier:parseval}
\begin{enumerate}
  \item Bila $f$ \omterm{def:b2:metric:continuity}{kontinu}, berkala-$2\pi$ dan $C^1$ sepotong-sepotong,
        maka $c_n(f') = \iu n\,c_n(f)$, deret Fourier $f$
        konvergen \emph{normal} pada $\R$, dan jumlahnya adalah $f$.
  \item (Parseval) Untuk setiap fungsi berkala-$2\pi$ yang \omterm{def:b2:metric:continuity}{kontinu}
        sepotong-sepotong, sebut saja $f$:\index{kesamaan Parseval}
        \[
        \frac{1}{2\pi}\int_{-\pi}^{\pi}\abs{f}^2
        = \sum_{n=-\infty}^{\infty} \abs{c_n(f)}^2
        = \frac{\abs{a_0}^2}{4} + \frac12\sum_{n\geq1}
        \bigl(\abs{a_n}^2 + \abs{b_n}^2\bigr).
        \]
        \emph{(Dibuktikan di sini untuk $f$ \omterm{def:b2:metric:continuity}{kontinu} yang $C^1$
        sepotong-sepotong; dan diterima tanpa bukti secara umum.)}
\end{enumerate}
\end{theorem}

\begin{proof}
(1) Pengintegralan parsial pada tiap keping $C^1$-nya (dengan suku
perbatasannya saling meniadakan berkat kekontinuan dan kekalaannya) memberi
$c_n(f') = \iu n c_n(f)$. Lalu, lewat Cauchy--Schwarz pada kedua keluarga
yang \omterm{def:b2:series:summable}{terjumlahkan} kuadratnya (\cref{prop:b2:fourier:bessel} bagi $f'$):
\[
\sum_{n \neq 0} \abs{c_n(f)} = \sum_{n\neq0}
\frac{\abs{c_n(f')}}{\abs n}
\leq \Bigl(\sum \abs{c_n(f')}^2\Bigr)^{1/2}
\Bigl(\sum_{n\neq0}\frac{1}{n^2}\Bigr)^{1/2} < \infty :
\]
jadi deret Fouriernya konvergen normal. Jumlahnya \omterm{def:b2:metric:continuity}{kontinu} dan
berimpit dengan $f$ di setiap titik menurut Dirichlet
(\cref{thm:b2:fourier:dirichlet}: sebab $f$ \omterm{def:b2:metric:continuity}{kontinu}): jadi deretnya
konvergen ke $f$ secara \omterm{def:b2:funcseq:def}{seragam}.

(2) Untuk $f$ yang demikian: $S_N f \to f$ secara \omterm{def:b2:funcseq:def}{seragam}, jadi $\norm{f - S_Nf}_2
\leq \norm{f - S_Nf}_\infty \to 0$, lalu Pythagoras
(sebab $\norm f_2^2 = \sum_{\abs n \leq N}\abs{c_n}^2 + \norm{f -
S_Nf}_2^2$) melewatkan limitnya. Adapun bentuk realnya sekadar tata buku dengan
$a_n, b_n$.
\end{proof}

\begin{example}[Batas ekor $C^1$, dibuat kuantitatif]\label{ex:b2:fourier:tailbound}
Bukti \cref{thm:b2:fourier:parseval}~(1) menyembunyikan sebuah taksiran
yang berguna. Untuk fungsi \omterm{def:b2:metric:continuity}{kontinu} yang $C^1$ sepotong-sepotong, sebut $f$,
Cauchy--Schwarz yang sama, diterapkan pada ekornya saja, memberi
\[
\sum_{\abs n > N}\abs{c_n(f)}
= \sum_{\abs n > N}\frac{\abs{c_n(f')}}{\abs n}
\leq \Bigl(\sum_{\abs n > N}\abs{c_n(f')}^2\Bigr)^{\!1/2}
\Bigl(\sum_{\abs n>N}\frac{1}{n^2}\Bigr)^{\!1/2}
\leq \norm{f'}_2\,\sqrt{\frac{2}{N}} ,
\]
dengan memakai Bessel bagi $f'$ dan $\sum_{n>N}n^{-2} \leq \frac1N$. Jadi
galat \omterm{def:b2:funcseq:def}{seragam} jumlah parsialnya menaati
\[
\norm{f - S_Nf}_\infty
\leq \sum_{\abs n>N}\abs{c_n(f)}
\leq \norm{f'}_2\,\sqrt{\frac2N} .
\]
Untuk $f(t) = \abs t$: berlaku $\norm{f'}_2 = 1$ (sebab turunannya
$\pm1$), jadi sepuluh suku saja sudah membangun ulang $\abs t$ secara seragam
dalam $\sqrt{0.2} \approx 0.45$, dan $N = 10^4$ dalam
$0.015$. Pelajaran penutupnya: satu turunan membeli laju
\omterm{def:b2:funcseq:def}{seragam} $\frac{1}{\sqrt N}$; dan bila dibandingkan dengan
dunia berkoefisien $\frac1n$ milik gelombang persegi (yang sama sekali tak
konvergen seragam), kamus
\cref{ex:b2:fourier:decaytable} memperoleh angkanya.
\end{example}

\begin{example}[Basel dan kawan-kawan]\label{ex:b2:fourier:basel}
Misalkan $f(t) = t$ pada $\intoo{-\pi}{\pi}$, yang diperluas secara
berkala-$2\pi$ (yakni gigi gergaji yang $C^1$ sepotong-sepotong). Setelah dihitung, $a_n = 0$ (berkat keganjilannya) dan
\[
b_n = \frac{1}{\pi}\int_{-\pi}^{\pi} t\sin nt\,\dd t
= \frac{2(-1)^{n+1}}{n} .
\]
Dirichlet di $t = \frac\pi2$ memulihkan $\frac\pi4 = 1 -
\frac13 + \frac15 - \dots$ milik Leibniz; sedangkan Parseval memberi
\[
\frac{1}{2\pi}\int_{-\pi}^{\pi} t^2\,\dd t = \frac{\pi^2}{3}
= \frac12\sum_{n\geq1}\frac{4}{n^2}
\quad\Longrightarrow\quad
\boxed{\;\sum_{n\geq1}\frac{1}{n^2} = \frac{\pi^2}{6}\;}
\]
--- yakni jumlah Basel milik Euler, dalam dua baris. Fungsi $f(t) = t^2$
serupa itu menghasilkan $\sum \frac1{n^4} = \frac{\pi^4}{90}$
(\cref{exo:b2:fourier:3}).
\end{example}

\begin{example}[Sebuah uraian lengkap dengan pemeriksaan bawaan: $\abs{\sin t}$]\label{ex:b2:fourier:abssin}
Fungsi $f(t) = \abs{\sin t}$ bersifat \omterm{def:b2:metric:continuity}{kontinu}, genap,
berkala-$\pi$ (sehingga berkala-$2\pi$), dan $C^1$ sepotong-sepotong.
Kegenapannya membunuh $b_n$; lalu $a_0 = \frac1\pi\int_{-\pi}^{\pi}
\abs{\sin t}\dd t = \frac4\pi$; dan untuk $n \geq 1$,
hasil-kali-ke-jumlah memberi
\[
a_n = \frac2\pi\int_0^\pi \sin t\cos nt\,\dd t
= \frac{1}{\pi}\int_0^\pi\bigl(\sin(1+n)t +
\sin(1-n)t\bigr)\dd t
= \frac2\pi\cdot\frac{1 + \cos n\pi}{1 - n^2}
\]
untuk $n \neq 1$ (dan $a_1 = 0$ secara langsung): jadi nol untuk $n$ ganjil, dan
$a_{2k} = \frac{-4}{\pi(4k^2-1)}$. Menurut
\cref{thm:b2:fourier:parseval}~(1) kekonvergenannya normal,
dan
\[
\abs{\sin t} = \frac{2}{\pi} - \frac{4}{\pi}
\sum_{k\geq1}\frac{\cos(2kt)}{4k^2 - 1}
\qquad (t \in \R) .
\]
Pemeriksaan bawaannya di $t = 0$: kesamaannya menuntut
$\sum_{k\geq1}\frac{1}{4k^2-1} = \frac12$, yang ditegaskan
teleskopnya:
\[
\sum_{k\geq1}\frac{1}{4k^2-1}
= \frac12\sum_{k\geq1}\Bigl(\frac{1}{2k-1} -
\frac{1}{2k+1}\Bigr) = \frac12 . \checkmark
\]
Pelajaran penutupnya: \omterm{def:b2:reduction:eigen}{spektrum} $\abs{\sin}$ hidup pada
frekuensi \emph{genap} belaka --- sebab menyearahkan sebuah sinus melipatduakan
kandungan frekuensinya, dan itulah sebabnya penyearah gelombang penuh berdengung pada
$100$ atau $120$ hertz, dua kali frekuensi jala-jalanya.
\end{example}

\begin{example}[Parseval sebagai alat hitung]\label{ex:b2:fourier:zetafourodd}
Parseval mengubah tiap uraian yang ada menjadi deret numerik secara borongan.
Terapkanlah ia pada $f(t) = \abs t$ (\cref{exo:b2:fourier:2}: dengan $a_0 =
\pi$, $a_n = \frac{-4}{\pi n^2}$ untuk $n$ ganjil, dan sisanya nol):
\[
\frac{1}{2\pi}\int_{-\pi}^{\pi}t^2\,\dd t = \frac{\pi^2}{3}
= \frac{a_0^2}{4} + \frac12\sum_{n \text{ ganjil}} a_n^2
= \frac{\pi^2}{4} +
\frac{8}{\pi^2}\sum_{n\text{ ganjil}}\frac{1}{n^4} ,
\]
sehingga
\[
\sum_{n\text{ ganjil}}\frac{1}{n^4}
= \frac{\pi^2}{8}\Bigl(\frac{\pi^2}{3} -
\frac{\pi^2}{4}\Bigr) = \frac{\pi^4}{96} .
\]
Periksa silang terhadap \cref{exo:b2:fourier:3}: memecah
$\sum\frac1{n^4}$ menjadi bagian ganjil dan genap memberi tata buku
bergaya-$\zeta$ $S = S_{\mathrm{ganjil}} + \frac{S}{16}$, jadi $S =
\frac{16}{15}\cdot\frac{\pi^4}{96} = \frac{\pi^4}{90}$ ---
persis nilai yang ditemukan di sana lewat fungsi yang berbeda. Dua
uraian, satu bilangan: jadi kekonsistenannya adalah isometri Parseval
yang bekerja. Pelajaran penutupnya: tiap uraian Fourier yang baru adalah
mesin pembangkit bagi kesamaan deret; dan Bagian II soal akhir
pekan menjelaskan mengapa mesinnya tak pernah dapat
membantah dirinya sendiri.
\end{example}

\begin{omfigure}
\begin{tikzpicture}
\begin{axis}[omaxis, width=9.2cm, height=5.6cm,
    xmin=-3.6, xmax=3.6, ymin=-1.5, ymax=1.5,
    xtick={-3.1416, 0, 3.1416},
    xticklabels={$-\pi$, $0$, $\pi$}, ytick={-1,1}]
  \addplot[omExo, thick] coordinates {(-3.1416,-1) (0,-1)};
  \addplot[omExo, thick] coordinates {(0,1) (3.1416,1)};
  \addplot[omDef, thick, domain=-3.3:3.3, samples=200]
    {4/pi*sin(deg(x))};
  \addplot[omThm, very thick, domain=-3.3:3.3, samples=400]
    {4/pi*(sin(deg(x)) + sin(deg(3*x))/3 + sin(deg(5*x))/5 +
     sin(deg(7*x))/7 + sin(deg(9*x))/9)};
\end{axis}
\end{tikzpicture}

{\small Gelombang persegi (kelabu) beserta \omterm{def:b2:fourier:coefficients}{jumlah parsial Fourier} $S_1$
(biru) dan $S_9$ (merah): jadi konvergen di setiap titik kekontinuannya, tetapi
dengan lonjakan $\sim 9\%$ yang bertahan di dekat lompatannya --- yakni
\emph{gejala Gibbs}\index{gejala Gibbs}. \omterm{def:b2:funcseq:def}{Kekonvergenan seragamnya}
gagal persis karena limitnya tak \omterm{def:b2:metric:continuity}{kontinu}.}
\end{omfigure}

\begin{example}[Penggeseran dan pemodulasian]\label{ex:b2:fourier:shiftrules}
Dua aturan satu baris membangkitkan banyak uraian dari satu. Untuk $a
\in \R$, setelah menyulihkan $s = t - a$:
\[
c_n\bigl(f(\cdot - a)\bigr)
= \frac{1}{2\pi}\int_{-\pi}^{\pi}f(t - a)\eu^{-\iu nt}\dd t
= \eu^{-\iu na}\,c_n(f)
\qquad\text{(penggeseran memodulasi spektrumnya)},
\]
lalu langsung dari definisinya,
\[
c_n\bigl(\eu^{\iu kt}f\bigr) = c_{n-k}(f)
\qquad\text{(pemodulasian menggeser spektrumnya)}.
\]
Contoh yang dikerjakan: gigi gergaji yang digeser sebesar $\pi$, yakni $g(t) = f(t -
\pi)$ dengan $f(t) = t$, punya koefisien-$b_n$
$(-1)^n\cdot\frac{2(-1)^{n+1}}{n} = -\frac2n$: jadi uraiannya
$g \sim -2\sum\frac{\sin nt}{n}$ bagi gigi gergaji yang melompat di
$0$ alih-alih di $\pi$ --- tanpa satu integral pun dihitung ulang. Pelajaran
penutupnya: penggeseran waktu hanya memutar fasanya, tak pernah amplitudonya
(sebab $\abs{c_n}$ invarian terhadap geseran), dan itulah sebabnya energi
(Parseval) dan kelas kekonvergenannya merupakan sifat \emph{bentuk}
sinyalnya, bukan sifat tempat jamnya dimulai.
\end{example}

\begin{remark}[Jebakan yang sering muncul]
\emph{(i) Tiga kekonvergenan, tiga mata uang:} \omterm{def:b2:funcseq:def}{titik demi titik}
(yakni Dirichlet: menuntut $C^1$ sepotong-sepotong, dan membayar \emph{titik tengahnya} di
tiap lompatan --- tak pernah nilai sepihaknya), \omterm{def:b2:funcseq:def}{seragam} (menuntut
limit yang \omterm{def:b2:metric:continuity}{kontinu}; jadi mustahil menyeberangi lompatan, dan Gibbs adalah
gejalanya yang kasatmata), dan rerata kuadratik (yakni Parseval: yang paling
tegar, dan buta terhadap titik individual). Sebutkanlah selalu yang mana
yang kauklaim. \emph{(ii) Tak ada penurunan suku demi suku secara
baku:} menurunkan deret gigi gergaji pada
\cref{ex:b2:fourier:basel} suku demi suku menghasilkan $\sum
2(-1)^{n+1}\cos nt$, yang sukunya bahkan tak menuju $0$ ---
sebab teorema pemindahan pada bab barisan fungsi menuntut
\omterm{def:b2:funcseq:def}{kekonvergenan seragam} deret \emph{turunannya}, dan lompatannya
menghancurkannya. Jadi kemulusan dulu, penurunan kemudian
(\cref{exo:b2:fourier:6} adalah kamusnya). \emph{(iii)
Jumlah parsial yang simetris:} teorema Dirichlet berbicara tentang $S_N =
\sum_{-N}^{N}$; jadi menata ulang atau menjumlahkan satu sisinya lebih dulu dapat
mengubah kedivergenan menjadi kekonvergenan dan sebaliknya. \emph{(iv)
Pergeseran penormalan:} konvensinya berbeda antarbuku
(dengan $\frac{1}{2\pi}$ atau $\frac1\pi$ di depan, berkala $2\pi$ atau
$1$); jadi invarian yang andal adalah hubungan keortonormalannya
--- maka hitunglah ulang $\langle e_m, e_n\rangle$ dalam konvensi yang
dipakai sebelum memercayai rumus apa pun.
\end{remark}

\begin{remark}[Di mana ini dipakai]
Parseval adalah benih teori $L^2$ bagi deret Fourier: sebab
jilid Tahun ke-3 merampungkan gambarannya (yakni eksponensialnya menjadi
basis Hilbert bagi $L^2$, dan pemetaan $f \mapsto (c_n)$ menjadi
isometri yang bijektif). Di dalam jilid ini, soal akhir pekan
membuktikan \emph{teorema Fejér} --- yakni rata-rata Cesàro deret
Fouriernya konvergen \omterm{def:b2:funcseq:def}{seragam} bagi setiap $f$ berkala yang \omterm{def:b2:metric:continuity}{kontinu}
--- yang meningkatkan Parseval umum yang diterima tanpa bukti menjadi
teorema, menghasilkan teorema Weierstrass trigonometri, lalu membayar
dua panen yang spektakuler: teorema pemerataan Weyl dan
ketaksamaan isoperimetrik. Adapun matematika terapan membaca bab
ini sehari-hari: \omterm{def:b2:reduction:eigen}{spektrum} sinyal, harmonik sistem yang bergetar,
dan transformasi Fourier cepat (dengan awatara hingganya tadi
\cref{exo:b2:hermitian:10}).
\end{remark}

\begin{remark}[Pandangan ke depan di dalam jilid ini]
Tiga bab bercakap-cakap dengan bab ini. Ke belakang: bab
\omterm{def:b2:hermitian:adjoint}{Hermitian} memasok geometrinya (yakni keluarga ortonormal,
proyeksi, Bessel), dan bab barisan fungsi memasok
analisisnya (yakni \omterm{def:b2:funcseq:def}{kekonvergenan seragam}, teorema pemindahan, identitas
hampiran --- sebab kernel Fejér bagi deret Fourier itu seperti
\omterm{thm:b2:funcseq:weierstrass}{polinomial Bernstein} bagi Weierstrass). Ke samping:
teori perbatasan bab deret pangkat kembali lewat \omterm{thm:b2:series:abel}{penjumlahan
Abel}, yang di sini diwujudkan kernel Poisson
(\cref{exo:b2:fourier:12}) --- dengan jari-jari cakramnya $r$ memainkan
peran parameter penjumlahannya. Ke depan: bab
persamaan diferensial menguraikan paksaan berkala
menjadi harmonik lalu menyuapkan tiap-tiapnya ke tanggapan frekuensi
osilatornya; dan resonansi terjadi ketika sebuah ragam Fourier masukannya
cocok dengan frekuensi alaminya, dan itulah sebabnya soal akhir pekannya
dan soal bab ini merupakan dua paruh satu cerita.
\end{remark}

\begin{omfigure}
\begin{tikzpicture}
\begin{axis}[omaxis, width=9.2cm, height=5.6cm,
    xmin=-3.6, xmax=3.6, ymin=-2.5, ymax=9.5,
    xtick={-3.1416, 0, 3.1416},
    xticklabels={$-\pi$, $0$, $\pi$}, ytick={8}]
  \addplot[omDef, thick, domain=0.02:3.3, samples=300]
    {sin(deg(8.5*x))/sin(deg(0.5*x))};
  \addplot[omDef, thick, domain=-3.3:-0.02, samples=300]
    {sin(deg(8.5*x))/sin(deg(0.5*x))};
  \addplot[omThm, very thick, domain=0.02:3.3, samples=300]
    {sin(deg(4*x))^2/(8*sin(deg(0.5*x))^2)};
  \addplot[omThm, very thick, domain=-3.3:-0.02, samples=300]
    {sin(deg(4*x))^2/(8*sin(deg(0.5*x))^2)};
  \node[font=\small, omDef] at (axis cs:1.05,5.2) {$D_8$};
  \node[font=\small, omThm] at (axis cs:-0.75,6.2) {$F_8$};
\end{axis}
\end{tikzpicture}

{\small \omterm{lem:b2:fourier:kernel}{Kernel Dirichlet} $D_8$ (biru) berayun dan mengambil
nilai negatif; sedangkan kernel Fejér $F_8$ (merah) bersifat taknegatif,
memusat di $0$, dan bererata $1$: yakni sebuah \emph{identitas
hampiran}. Kepositifan itulah persis yang tak dimiliki \omterm{lem:b2:fourier:kernel}{kernel Dirichlet},
dan yang membuat teorema Fejér pada soal akhir pekan
menjadi tanpa syarat.}
\end{omfigure}

\section{Latihan}

\begin{exercise}[$\star$]\label{exo:b2:fourier:1}
Hitunglah \omterm{def:b2:fourier:coefficients}{koefisien Fourier} gelombang persegi (dengan $f = -1$ pada
$\intoo{-\pi}{0}$ dan $+1$ pada $\intoo{0}{\pi}$), nyatakan kesimpulan
Dirichlet di $t = \frac\pi2$ dan di lompatan $t = 0$, lalu peroleh kembali
deret Leibniz.
\end{exercise}

\begin{exercise}[$\star$]\label{exo:b2:fourier:2}
Uraikan $f(t) = \abs t$ (dengan $\abs t \leq \pi$, berkala-$2\pi$) menjadi
deret Fourier; benarkan kekonvergenan \emph{normal}nya; lalu nilailah di $t =
0$ untuk memperoleh $\sum_{k\geq0}\frac{1}{(2k+1)^2} = \frac{\pi^2}{8}$,
kemudian turunkan ulang Basel darinya.
\end{exercise}

\begin{exercise}[$\star$]\label{exo:b2:fourier:3}
Uraikan $f(t) = t^2$ (dengan $\abs t \leq \pi$) lalu turunkan
\[
\sum_{n\geq1}\frac{(-1)^{n+1}}{n^2} = \frac{\pi^2}{12},
\qquad
\sum_{n\geq1}\frac{1}{n^4} = \frac{\pi^4}{90}
\quad\text{(Parseval)}.
\]
\end{exercise}

\begin{exercise}[$\star\star$]\label{exo:b2:fourier:4}
Misalkan $f$ \omterm{def:b2:metric:continuity}{kontinu} berkala-$2\pi$ dengan $c_n(f) = 0$ untuk setiap
$n$. Buktikan $f = 0$ \emph{(lewat Parseval --- untuk kelas mana ia dibuktikan
di sini? benarkanlah bahwa kekontinuan ditambah $C^1$ sepotong-sepotong dapat dilepas
dengan menerima Parseval umum, atau berikanlah hujah kepadatannya secara garis besar)}.
\end{exercise}

\begin{exercise}[$\star\star$]\label{exo:b2:fourier:5}
Untuk $\alpha \notin \Z$, uraikan $f(t) = \cos(\alpha t)$ (dengan $\abs t
\leq \pi$) lalu turunkan uraian pecahan parsial bagi
kotangennya:
\[
\pi\cot(\pi\alpha) = \frac{1}{\alpha} + \sum_{n\geq1}
\frac{2\alpha}{\alpha^2 - n^2} .
\]
\end{exercise}

\begin{exercise}[$\star\star$]\label{exo:b2:fourier:6}
Buktikan bahwa bila $f$ berkala-$2\pi$ dan $C^k$ dengan $f^{(k)}$
\omterm{def:b2:metric:continuity}{kontinu} sepotong-sepotong, maka $c_n(f) = O\bigl(\abs n^{-k}\bigr)$: jadi
kemulusan sinyalnya $=$ peluruhan \omterm{def:b2:reduction:eigen}{spektrumnya}.
\end{exercise}

\begin{exercise}[$\star\star\star$]\label{exo:b2:fourier:7}
(Ketaksamaan Wirtinger) Misalkan $f$ berkelas $C^1$, berkala-$2\pi$, dengan
$\int_{-\pi}^{\pi} f = 0$. Buktikan
\[
\int_{-\pi}^{\pi} \abs{f}^2 \leq \int_{-\pi}^{\pi} \abs{f'}^2 ,
\]
dengan kesamaan bila dan hanya bila $f(t) = a\cos t + b\sin t$. \emph{(Lewat Parseval pada
kedua ruasnya; lalu bandingkan $\abs{c_n}^2$ dan $n^2\abs{c_n}^2$.)}
\end{exercise}

\begin{exercise}[$\star\star\star$]\label{exo:b2:fourier:8}
(Konstanta Gibbs) Untuk gelombang persegi pada
\cref{exo:b2:fourier:1}, nilailah jumlah parsialnya di $x_N =
\frac{\pi}{2N}$: dengan menulis $u_k = \frac{(2k+1)\pi}{2N}$ dan $\Delta u
= \frac{\pi}{N}$, tunjukkanlah bahwa
\[
S_{2N-1}\Bigl(\frac{\pi}{2N}\Bigr)
= \frac{4}{\pi}\sum_{k=0}^{N-1} \frac{\sin u_k}{2k+1}
= \frac{2}{\pi}\sum_{k=0}^{N-1} \frac{\sin u_k}{u_k}\,\Delta u
\xrightarrow[N\to\infty]{}
\frac{2}{\pi}\int_0^{\pi}\frac{\sin u}{u}\,\dd u \approx 1.179 :
\]
yakni sebuah jumlah Riemann bagi $\frac{2}{\pi}\cdot\frac{\sin u}{u}$ pada
$\intcc{0}{\pi}$ di titik tengahnya. Simpulkan bahwa lonjakan di luar
nilai lompatannya $1$ tidak lenyap ketika $N \to \infty$.
\end{exercise}

\begin{exercise}[$\star$]\label{exo:b2:fourier:9}
Uraikan $\cos^3 t$ dan $\sin^2 t\,\cos t$ menjadi deret Fourier
\emph{(linearkan; sebab sebuah polinomial trigonometri adalah deret Fouriernya
sendiri, berkat ketunggalan koefisiennya)}. Berapakah $c_n$,
$a_n$, dan $b_n$ untuk masing-masingnya?
\end{exercise}

\begin{exercise}[$\star\star$]\label{exo:b2:fourier:10}
Misalkan $a > 0$ dan $f(t) = \eu^{at}$ pada $\intoc{-\pi}{\pi}$, yang
diperluas secara berkala-$2\pi$. Hitunglah
\[
c_n(f) = \frac{(-1)^n\sinh(a\pi)}{\pi(a - \iu n)} ,
\]
lalu terapkan teorema Dirichlet di lompatan $t = \pi$, kemudian turunkan
uraian pecahan parsial bagi kotangen hiperboliknya:
\[
\coth(\pi a) = \frac{1}{\pi a} + \sum_{n\geq1}
\frac{2a}{\pi(a^2 + n^2)} .
\]
\end{exercise}

\begin{exercise}[$\star\star$]\label{exo:b2:fourier:11}
(Konvolusi) Untuk $f, g$ yang \omterm{def:b2:metric:continuity}{kontinu} dan berkala-$2\pi$, definisikan
\[
(f * g)(x) = \frac{1}{2\pi}\int_{-\pi}^{\pi} f(x - t)\,g(t)\,
\dd t .
\]
Tunjukkan bahwa $f * g = g * f$, bahwa $c_n(f*g) = c_n(f)\,c_n(g)$
\emph{(untuk menukarkan kedua integral integran yang \omterm{def:b2:metric:continuity}{kontinu},
bandingkanlah kedua fungsi batas atasnya: sebab keduanya lenyap di
ujung kirinya dan punya turunan yang sama, berkat kekontinuan dan
penurunan di bawah tanda integralnya)}, dan bahwa $S_N(f) =
f * D_N$ untuk \omterm{lem:b2:fourier:kernel}{kernel Dirichletnya}. (Rata-rata Fejér pada soal akhir
pekan juga berupa konvolusi, yakni $\sigma_N(f) = f *
F_N$.)
\end{exercise}

\begin{exercise}[$\star\star\star$]\label{exo:b2:fourier:12}
(Kernel Poisson: rata-rata Abel deret Fourier) Untuk $0 \leq r <
1$ tetapkan $P_r(t) = \sum_{n\in\Z} r^{\abs n}\eu^{\iu nt}$.
\begin{enumerate}
  \item Jumlahkan kedua deret geometrinya lalu tunjukkan
        \[
        P_r(t) = \frac{1 - r^2}{1 - 2r\cos t + r^2} > 0,
        \qquad
        \frac{1}{2\pi}\int_{-\pi}^{\pi}P_r = 1 .
        \]
  \item Tunjukkan bahwa untuk $\delta \leq \abs t \leq \pi$: $P_r(t)
        \leq \frac{1 - r^2}{1 - 2r\cos\delta + r^2} \to 0$ ketika
        $r \to 1^-$, secara \omterm{def:b2:funcseq:def}{seragam}.
  \item Turunkan bahwa untuk setiap fungsi \omterm{def:b2:metric:continuity}{kontinu} berkala-$2\pi$, sebut $f$,
        \emph{rata-rata Abel}nya $(f * P_r)(x) = \sum_n
        r^{\abs n}c_n(f)\,\eu^{\iu nx}$ konvergen ke $f$
        secara \omterm{def:b2:funcseq:def}{seragam} ketika $r \to 1^-$ --- yakni saudara
        berparameter \omterm{def:b2:metric:continuity}{kontinu} bagi teorema Fejér, sekaligus penjelmaan
        Fourier bagi \omterm{thm:b2:series:abel}{penjumlahan Abel} dari bab deret
        pangkat.
\end{enumerate}
\end{exercise}

\section{Soal: Teorema Fejér dan panennya}

\begin{problem}\label{pb:b2:fourier:1}
Teorema Dirichlet menuntut $f$ yang $C^1$ sepotong-sepotong; sedangkan untuk $f$
yang sekadar \omterm{def:b2:metric:continuity}{kontinu} jumlah parsial $S_N(f)$ dapat berulah. Adapun temuan
Fejér: \emph{rata-rata Cesàro}nya tak pernah berulah. Mesinnya adalah
kepositifan kernel Fejér, dan panennya
melimpah: hampiran trigonometri yang \omterm{def:b2:funcseq:def}{seragam} (Weierstrass),
ketunggalan \omterm{def:b2:fourier:coefficients}{koefisien Fourier}, Parseval bagi setiap fungsi
bab ini (sehingga membuang ``yang diterima tanpa bukti'' pada
\cref{thm:b2:fourier:parseval}), teorema pemerataan Weyl,
dan --- sebagai mahkota satu abad geometri --- 
ketaksamaan isoperimetrik.\index{teorema Fejer@teorema Fejér}\index{pemerataan}\index{ketaksamaan isoperimetrik}
Di sepanjang soal ini, $f$ berkala-$2\pi$ dan \omterm{def:b2:metric:continuity}{kontinu} sepotong-sepotong,
lalu
\[
\sigma_N(f) = \frac{S_0(f) + S_1(f) + \dots +
S_{N-1}(f)}{N} .
\]

\medskip
\noindent\textbf{Bagian I --- Kernel Fejér.}
\begin{enumerate}
  \item Tunjukkan bahwa $\sigma_N(f)(x) =
        \frac{1}{2\pi}\int_{-\pi}^{\pi}f(x+u)\,F_N(u)\,\dd u$
        dengan $F_N = \frac{D_0 + \dots + D_{N-1}}{N}$, dan
        bahwa $\frac{1}{2\pi}\int_{-\pi}^{\pi}F_N = 1$.
  \item Buktikan bentuk tertutupnya, untuk $u \notin 2\pi\Z$:
        \[
        F_N(u) = \frac{1}{N}\,
        \frac{\sin^2\bigl(\frac{Nu}{2}\bigr)}
        {\sin^2\bigl(\frac u2\bigr)} \;\geq\; 0
        \]
        \emph{(jumlahkan $\sin\bigl((n+\frac12)u\bigr)$ sebagai
        bagian imajiner sebuah deret geometri)}.
  \item Tunjukkan taksiran pemusatannya: untuk $0 < \delta \leq
        \abs u \leq \pi$,
        \[
        F_N(u) \leq \frac{1}{N\sin^2\frac\delta2}
        \xrightarrow[N\to\infty]{} 0
        \quad\text{secara seragam} :
        \]
        jadi $(F_N)$ merupakan identitas hampiran yang positif.
  \item (Teorema Fejér) Buktikan: bila $f$ \omterm{def:b2:metric:continuity}{kontinu} dan
        berkala-$2\pi$, maka $\sigma_N(f) \to f$
        \emph{secara \omterm{def:b2:funcseq:def}{seragam}} pada $\R$ \emph{(pecahlah integral
        $\bigl(f(x+u) - f(x)\bigr)F_N(u)$ di $\abs u =
        \delta$; lalu pakailah Heine beserta pertanyaan 1--3)}.
  \item Untuk $f$ yang \omterm{def:b2:metric:continuity}{kontinu} sepotong-sepotong, tunjukkan versi
        \omterm{def:b2:funcseq:def}{titik demi titiknya} $\sigma_N(f)(x) \to \frac{f(x^+) +
        f(x^-)}{2}$ di setiap $x$, beserta batas \omterm{def:b2:funcseq:def}{seragamnya}
        $\norm{\sigma_N(f)}_\infty \leq \norm f_\infty$
        \emph{(berkat kepositifannya!)}.
\end{enumerate}

\medskip
\noindent\textbf{Bagian II --- Weierstrass, ketunggalan,
Parseval.}
\begin{enumerate}[resume]
  \item (Weierstrass trigonometri) Turunkan: bahwa setiap fungsi
        \omterm{def:b2:metric:continuity}{kontinu} berkala-$2\pi$ merupakan limit \omterm{def:b2:funcseq:def}{seragam}
        polinomial trigonometri.
  \item (Ketunggalan) Turunkan: bahwa $f$ \omterm{def:b2:metric:continuity}{kontinu} dengan $c_n(f) = 0$
        untuk setiap $n$ pastilah nol secara identik --- jadi dua fungsi
        berkala yang \omterm{def:b2:metric:continuity}{kontinu} dengan \omterm{def:b2:fourier:coefficients}{koefisien Fourier} yang sama
        berimpit (\cref{exo:b2:fourier:4}, kini tanpa
        penerimaan tanpa bukti).
  \item (Parseval, kasus \omterm{def:b2:metric:continuity}{kontinu}) Dengan memakai sifat proyeksi
        milik $S_N$ (\cref{prop:b2:fourier:bessel}) dan
        $\sigma_Nf \in \mathcal T_{N-1} \subseteq \mathcal
        T_N$, buktikanlah
        \[
        \norm{f - S_Nf}_2 \leq \norm{f - \sigma_Nf}_2
        \leq \norm{f - \sigma_Nf}_\infty
        \xrightarrow[N\to\infty]{} 0 ,
        \]
        lalu simpulkan kesamaan Parseval bagi setiap fungsi
        \emph{\omterm{def:b2:metric:continuity}{kontinu}} yang berkala-$2\pi$, sebut $f$.
  \item (Parseval, kasus \omterm{def:b2:metric:continuity}{kontinu} sepotong-sepotong) Diberikan $f$
        yang \omterm{def:b2:metric:continuity}{kontinu} sepotong-sepotong dan $\varepsilon > 0$, bangunlah
        $g$ berkala yang \omterm{def:b2:metric:continuity}{kontinu} dengan $\norm{f - g}_2 \leq
        \varepsilon$ \emph{(gantilah $f$ dengan interpolasi
        afin pada interval mungil di sekitar lompatannya)},
        lalu turunkan $\norm{f - S_Nf}_2 \to 0$
        \emph{(pakailah Bessel: $\norm{S_Nh}_2 \leq \norm h_2$)}: jadi
        Parseval berlaku dalam keumuman penuh yang dinyatakan pada
        \cref{thm:b2:fourier:parseval} --- sehingga ``yang diterima tanpa bukti''
        telah lenyap.
  \item (Tak ada Gibbs bagi Fejér) Bandingkan dengan
        \cref{exo:b2:fourier:8}: tunjukkan bahwa untuk gelombang persegi
        $f$ berlaku $\abs{\sigma_N(f)} \leq 1$ di mana-mana, untuk setiap
        $N$ --- jadi perata-rataan Cesàro menghapus lonjakan yang
        menghantui $S_N$. Jelaskanlah dalam satu kalimat sifat mana
        milik $F_N$ yang bertanggung jawab.
\end{enumerate}

\medskip
\noindent\textbf{Bagian III --- Laju.}
\begin{enumerate}[resume]
  \item Buktikan kedua batas kernelnya, untuk $0 < \abs u \leq
        \pi$:
        \[
        F_N(u) \leq N,
        \qquad
        F_N(u) \leq \frac{\pi^2}{N u^2}
        \]
        \emph{(untuk yang pertama, $\abs{\sin N\theta} \leq
        N\abs{\sin\theta}$ secara induktif; untuk yang kedua,
        $\sin\frac u2 \geq \frac{u}{\pi}$ pada
        $\intcc{0}{\pi}$)}.
  \item Turunkan taksiran momen pertamanya
        \[
        \frac{1}{2\pi}\int_{-\pi}^{\pi}\abs u\,F_N(u)\,\dd u
        \;\leq\; \frac{C\,\ln N}{N}
        \qquad (N \geq 2)
        \]
        untuk sebuah konstanta yang gamblang \emph{(pecahlah di $\abs u =
        \frac1N$)}.
  \item Simpulkan: bahwa bila $f$ \omterm{def:b2:metric:continuity}{Lipschitz} berkonstanta $L$ dan berkala-$2\pi$,
        maka
        \[
        \norm{\sigma_N f - f}_\infty \leq \frac{C\,L\ln N}{N} .
        \]
  \item (Penjenuhan) Hitunglah $\sigma_N(e_1)$ untuk $e_1(t) =
        \eu^{\iu t}$ lalu tunjukkan $\norm{\sigma_Ne_1 - e_1}_\infty
        = \frac1N$: jadi bahkan untuk fungsi termulus sekalipun, Fejér
        konvergen tak lebih cepat daripada $\frac1N$ --- yakni analogi
        persis bagi penjenuhan Bernstein pada soal akhir pekan bab
        barisan fungsi.
  \item (Pelokalan) Tunjukkan: bahwa bila $f$ (yang \omterm{def:b2:metric:continuity}{kontinu} sepotong-sepotong)
        lenyap pada $\intoo{x - \delta}{x + \delta}$, maka
        $\sigma_N(f)(x) \to 0$, sebuas apa pun $f$ di tempat lain
        --- jadi kekonvergenan rata-ratanya di $x$ hanya melihat $f$ di dekat
        $x$.
\end{enumerate}

\medskip
\noindent\textbf{Bagian IV --- Teorema pemerataan Weyl.}
Sebuah barisan $(x_n)_{n\geq1}$ di $\intco{0}{1}$ disebut
\emph{terdistribusi merata} apabila, untuk setiap interval
$\intcc{a}{b} \subseteq \intcc{0}{1}$,
\[
\frac{\#\{n \leq N : x_n \in \intcc ab\}}{N}
\xrightarrow[N\to\infty]{} b - a .
\]
\begin{enumerate}[resume]
  \item Tunjukkan bahwa $(x_n)$ terdistribusi merata segera setelah
        $\frac1N\sum_{n\leq N}f(x_n) \to \int_0^1 f$ untuk setiap fungsi
        \emph{\omterm{def:b2:metric:continuity}{kontinu}} yang berkala-$1$, sebut $f$ \emph{(jepitlah indikator
        $\intcc ab$ di antara dua fungsi \omterm{def:b2:metric:continuity}{kontinu} yang
        afin sepotong-sepotong dan integralnya berselisih
        $\varepsilon$)}.
  \item (Kriteria Weyl, kecukupannya) Andaikan
        \[
        \frac{1}{N}\sum_{n=1}^{N}\eu^{2\iu\pi kx_n}
        \xrightarrow[N\to\infty]{} 0
        \qquad\text{untuk setiap } k \in \Z\setminus\{0\} .
        \]
        Tunjukkan $\frac1N\sum f(x_n) \to \int_0^1f$ mula-mula untuk
        polinomial trigonometri, lalu untuk setiap fungsi \omterm{def:b2:metric:continuity}{kontinu}
        berkala-$1$, sebut $f$, lewat pertanyaan 6 (yang diangkut ke kala
        $1$): jadi bersama pertanyaan 16, $(x_n)$ terdistribusi merata.
  \item Misalkan $\alpha$ irasional dan $x_n = \{n\alpha\}$
        (yakni bagian pecahannya). Batasilah jumlah geometrinya
        \[
        \Bigl|\sum_{n=1}^{N}\eu^{2\iu\pi kn\alpha}\Bigr|
        \leq \frac{2}{\abs{1 - \eu^{2\iu\pi k\alpha}}}
        \qquad (k \neq 0),
        \]
        lalu simpulkan \emph{teorema Weyl}: bahwa $(\{n\alpha\})$
        terdistribusi merata di $\intco{0}{1}$.
  \item Turunkan bahwa $(\{n\alpha\})$ padat di
        $\intcc{0}{1}$ untuk $\alpha$ yang irasional, lalu jelaskanlah dalam
        satu kalimat mengapa pemerataan tegas lebih kuat
        daripada kepadatan.
  \item (Digit terdepan) Buktikan bahwa proporsi bilangan bulat
        $n \leq N$ yang membuat $2^n$ berdigit desimal terdepan
        $1$ menuju $\log_{10}2 \approx 0.301$ \emph{(sebab digit
        terdepannya $1$ berarti $\{n\log_{10}2\} \in
        \intco{0}{\log_{10}2}$; lalu tunjukkan $\log_{10}2$
        irasional)}.
\end{enumerate}

\medskip
\noindent\textbf{Bagian V --- Ketaksamaan isoperimetrik.} Misalkan
$\Gamma$ kurva $C^1$ tertutup sederhana berpanjang $L$ yang melingkupi
luas bertanda $A$, yang diparameterkan lewat panjang busur terskala: $z(t) =
x(t) + \iu y(t)$, berkala-$2\pi$, dengan $\abs{z'(t)} =
\frac{L}{2\pi}$ yang tetap; lalu luas yang dilingkupinya adalah
\[
A = \frac12\int_0^{2\pi}\bigl(x\,y' - y\,x'\bigr)\dd t
= \frac{1}{2}\,\Im\int_0^{2\pi}\conj{z}\,z'\,\dd t
\]
(yang di sini diambil sebagai definisi luas bertandanya; sedangkan bab
integral lipat membuktikan bahwa ia sepakat dengan yang intuitif,
lewat rumus Green).
\begin{enumerate}[resume]
  \item Uraikan $z(t) = \sum_{n\in\Z}c_n\eu^{\iu nt}$ (yakni
        deret sebuah fungsi $C^1$, yang konvergen normal) lalu
        buktikan, lewat Parseval yang diterapkan pada $z'$:
        \[
        \frac{L^2}{2\pi} = \int_0^{2\pi}\abs{z'}^2\dd t
        = 2\pi\sum_{n\in\Z}n^2\abs{c_n}^2 .
        \]
  \item Buktikan pula $A = \pi\sum_{n\in\Z}
        n\,\abs{c_n}^2$ \emph{(lewat Parseval dalam bentuk terpolarisasinya:
        $\frac{1}{2\pi}\int\conj f g = \sum
        \conj{c_n(f)}c_n(g)$, yang diterapkan pada $f = z$ dan $g = z'$)}.
  \item (Hurwitz) Simpulkan:
  \item (Hurwitz) Conclude:
        \[
        L^2 - 4\pi A = 4\pi^2\sum_{n\in\Z}
        (n^2 - n)\abs{c_n}^2 \;\geq\; 0 ,
        \]
        dengan kesamaan bila dan hanya bila $z(t) = c_0 +
        c_1\eu^{\iu t}$ --- yakni sebuah lingkaran. Itulah ketaksamaan
        isoperimetrik: di antara kurva tertutup berpanjang $L$, hanya
        lingkaran yang melingkupi luas $\frac{L^2}{4\pi}$.
  \item Pemeriksaan kewarasannya: periksalah kesamaannya bagi lingkaran
        berjari-jari $R$ dan ketaksamaan tegasnya bagi bujur sangkar
        bersisi $a$; lalu jelaskan mengapa $n^2 - n \geq 0$ untuk setiap
        bilangan bulat $n$, termasuk yang negatif, dan di mana
        kelajuan tetap parameterisasinya dipakai.
  \item Rangkuman. Satu kalimat untuk masing-masing: (i) satu
        sifat $F_N$ yang darinya Bagian I--III mengalir, dan
        yang tak dimiliki $D_N$; (ii) bagaimana penjumlahan Cesàro di sini
        berkaitan dengan soal akhir pekan bab deret pangkat
        (yakni Frobenius); (iii) panen mana yang hanya memakai
        Weierstrass (pertanyaan 6) dan mana yang menuntut
        Parseval penuh; (iv) satu kalimat tentang apa yang ditambahkan jilid
        Tahun ke-3 (yakni kelengkapan $L^2$: deret Fourier sebagai basis
        Hilbert).

\end{enumerate}
\end{problem}
